\documentclass[12pt,a4paper]{article}
\usepackage{ctex}
\usepackage{amsmath,amssymb}
\usepackage{geometry}
\geometry{left=2.5cm,right=2.5cm,top=2.5cm,bottom=2.5cm}
\usepackage{enumitem}
\usepackage{soul} % 导入 soul 包

\usepackage{color, xcolor} % 颜色包，color 必须导入，xcolor 建议导入
\usepackage{titlesec}
\usepackage{hyperref}

% 标题格式
\titleformat{\section}{\centering\bfseries\Large}{\thesection}{1em}{}
\titleformat{\subsection}{\bfseries\large}{\thesubsection}{1em}{}

\begin{document}

\begin{center}
\LARGE\bfseries 2026大学物理（上）期中复习安排
\end{center}

\colorbox{yellow}{黄色内容是想对你说的话}
\textcolor{blue}{蓝色字体是我的笔记}
\vspace{1em}

本次复习将采用问题导向结构，从运动学部分到流体力学部分，深入浅出地对大学物理复习进行指导，从易到难，简单考验计算能力。本次复习计划旨在令同学们深入学习后知其然，且知其所以然。本复习安排大概需要七天完成。希望同学们能够收获到自己满意的成绩！

\vspace{1em}

该部分是善用教材部分。请你再次阅读教材，跟随节奏进行研读。本复习计划在教材部分并没有答案，如果你有疑问，可以自行查阅资料或者问人工智能。本复习计划没有知识点整理，是因为这些基础部分需要在练习中自行感悟。本复习安排主要任务是开拓视野与寻求自我反思。

\section{一、运动学部分}
\begin{enumerate}[label=\arabic*.,leftmargin=*]
    \item 矢量的思维
    
    请用简短的中文解释
    \[
    \frac{d\boldsymbol{r}}{dt},\ \frac{d|\boldsymbol{r}|}{dt},\ \left|\frac{d\boldsymbol{r}}{dt}\right|,\ \left|\frac{d|\boldsymbol{r}|}{dt}\right|
    \]
    的物理意义。
    
    \textcolor{blue}{\colorbox{yellow}{与原点距离的变化率（？）}，速度，与原点距离的变化率的大小，速度的大小}

    \item 阅读例1-2.思考：速度分解到底是怎样分解的呢？
    
    \textcolor{blue}{沿任何两个你想要的方向分解（可以是能够简化问题的）}

    \item 阅读1.4曲线运动之前，请思考：为什么 $\boldsymbol{e}_x$ 对时间的微分为0，而 $\boldsymbol{e}_r$（极坐标系）微分就不是0了呢？自然坐标系情况又如何？
    
    \textcolor{blue}{$\boldsymbol{e}_x$ 是一段固定的向量，$\boldsymbol{e}_r$ 是可以随时间变换的。自然坐标系同样会变换，微分不是0。}

    \item 阅读1.4.1抛体运动。非常简单的，水平面的最远射程怎么计算？
    
    困难而不要求掌握，但是可以把结论搜出来记住的，如果是在有坡度的斜坡上的最远射程呢？
    
    其次，抛体运动是两种运动的叠加，并且哪怕有其他和速度相关的力（如空气阻力）的作用，也仍然是两种运动的叠加。
    
    尝试列出在有大小与速度大小成正比的，方向与速度反向的空气阻力的作用下的抛体运动的参数方程。
    
    \textcolor{blue}{
    \colorbox{yellow}{非常好的题目}\\
    列式后注意到$v_x$是变化的，可以求出$v_x (t)$，再次积分求出$x(t)$。
    不能直接套用$x=v_0*t+\frac{1}{2}*a*t^2$，加速度在变化。
    \[
    \begin{cases}
    x(t)=\dfrac{v_{x0}}{\gamma}\left(1-e^{-\gamma t}\right)\\[6pt]
    y(t)=\left(\dfrac{v_{y0}}{\gamma}+\dfrac{g}{\gamma^2}\right)\left(1-e^{-\gamma t}\right)-\dfrac{g}{\gamma}\,t
    \end{cases}
    \]
    其中 $\gamma=\dfrac{k}{m}$，$v_{x0}=v_0\cos\theta$，$v_{y0}=v_0\sin\theta$。
    }

    \item 阅读例1-4，考试可能会出一个这种形式的填空题。
    
    其次你可能无法在原表达式里找到时间$t$。
    
    对于这类式子的解构，有时候需要考虑以速度作为桥梁。在位移（路程）与加速度之间，速度是方便的一种关系。
    
    据他校同学提供的思路，可以尝试理解
    \[
    dt=\frac{dv}{a},\quad dt=\frac{dx}{v}
    \]

    \textcolor{blue}{
        \colorbox{yellow}{
        写作
    $a=\frac{dv}{dt},\quad v=\frac{dx}{dt}$
    也许会更好理解}\\
    依旧求矢量容易漏答方向}

    \item 习题大部分有做的需要。其中需要加以重视的：1-2（思考牵连速度）；1-5（考试可能出现）；1-10（做完可以研究其轨迹，发现本题妙处）；1-11与1-15（同类题，且有一定难度）
    
    \colorbox{yellow}{$v=\frac{H}{H-h}v_0$，哪里有什么牵连速度？}

    \colorbox{yellow}{是把人当作参考系吗？那是否可以问影子移动相对人的速度？}

    \colorbox{yellow}{$v_{\text{相对人}}=\frac{h}{H-h}v_0$?}
    
    \textcolor{blue}{1-5证明指向椭圆中心，除了叉乘=零向量，可以直接说明$\vec{a}=-\omega^2\vec{r}$}

    \textcolor{blue}{1-10表示变量，按需积分}

    \textcolor{blue}{1-11 圆盘参考系，引入惯性力来简化分析（用他们去等效真实的力）（圆盘参考系里用惯性力和真实的力可以列出平衡方程）\\
    惯性离心力=$m\omega^2 r$，向外；科里奥利力=$2m\vec{v}\times\vec{\omega} $，方向垂直于速度和旋转轴的平面内\\
    速度和加速度都分别求出径向和横向的\\
    }
\end{enumerate}

\section{二、牛顿运动定律}
\begin{enumerate}[label=\arabic*.,leftmargin=*]
    \item 量纲。阅读2.1，虽然不考，但应知应会。“量纲可以定性表示出物理量与基本量之间的关系，可以有效地应用它进行单位换算，可以用它来检查物理公式正确与否，还可以通过它来推知某些物理规律。”
    
    你可以保持做完题检查量纲的习惯，可以降低错误率。在式子中尽可能多地将量纲为1的式子分离出来，方便进行量纲检验。


    \item 阅读2.3.1，记住惯性系与非惯性系的定义。试图用自己的语言解释。\\
    \textcolor{blue}{惯性系：满足牛一；加速上升的电梯里若悬浮一个不受力（包括重力）的物体，此时它不受力但是运动状态改变，所以加速电梯是非惯性系}

    \item 阅读2.4，万有引力和重力是一个东西吗？区别在哪？\\
    \textcolor{blue}{引力可分解为重力和向心力，向心力指向地轴，其后确定重力}

    \item 阅读例2-3.发现对于相对复杂的系数，可以用比例系数表示。但是作为做题习惯，尽可能地在最后告诉改卷老师，“其中 $b=6\pi r\eta$。”\\
    \colorbox{yellow}{什么斯托克斯公式}


    \item （思维拓展）阅读例2-4.已知对于理想流体，比如这道题的液体，水面是等势面（但不一定单指重力势哦）。你有什么想法吗？在阅读完本复习计划的四（2）和四（3）后，询问人工智能直到你可以彻底理解其中的每一个点。此后，想想潮汐。\\
    \colorbox{yellow}{水面一定是等压面。关于重力势，倘若旋转起来，恐怕不是相等的}\\
    \colorbox{yellow}{又看了下，建议把表述改为“所有势能之和相等，不只有重力势”相关的}\\
    \textcolor{blue}{潮汐因为地月位置，引力不同，导致等势面-海面变形，形成潮汐}

    \item 阅读2.7，离心力是什么？\\
    \textcolor{blue}{惯性力}


    \item 阅读2.8，自主推导科里奥利力。你只用推导物体沿径向方向运动就好。它和习题1-11与1-15有异曲同工之妙。\\
    \textcolor{blue}{
$r=ut,\quad v_\theta=r\omega=ut\omega$\\
第一项（$r$ 增大）：
\[
a_1=\frac{\mathrm{d}v_\theta}{\mathrm{d}t}=u\omega
\]
第二项（径向速度方向转动）：
\[
\vec{v}_r=u\vec{e}_r,\quad
\frac{\mathrm{d}\vec{e}_r}{\mathrm{d}t}=\omega\vec{e}_\theta
\Rightarrow
a_2=\left|\frac{\mathrm{d}\vec{v}_r}{\mathrm{d}t}\right|=u\omega
\]
总切向加速度：
\[
a_\theta=a_1+a_2=2u\omega
\]
为什么考虑这两个？\\
因为速度本就由径向和切向两个矢量叠加而成，径向的方向在改变（受加速度），切向的速度在改变
}
    

    \item 习题需要加以重视的：2-11（很有意思，尽可能多解）；2-13（注意算对）；2-14（方程可能会列的比较多）；2-15（《难题集萃》第二章题1.结果是15圈，考不到这么难，有兴趣可以挑战）。\\
    \textcolor{blue}{2-14 绳子两端的加速度一致，列四个方程\\
    2-15
柔索绕固定圆柱的临界静摩擦公式（欧拉公式）：
\[
T_2 = T_1 e^{\mu \theta}
\]
其中：
- 拉力 \(T_2 = m_0 g,\ T_1 = m g\)，约去重力 \(g\) 得 \(\dfrac{m_0}{m} = e^{\mu\theta}\)；
- \(\theta = 2\pi n\) 为绳索包角，\(n\) 为缠绕圈数；
- 摩擦因数 \(\mu=0.05\)，\(m_0=1000\ \text{kg}\)，\(m=10\ \text{kg}\)。
对公式取自然对数：
\[
\ln\frac{m_0}{m} = \mu \cdot 2\pi n
\]
解得圈数：
\[
n = \frac{\ln(m_0/m)}{2\pi\mu}
= \frac{\ln100}{2\pi\times0.05}
\approx 14.66
\]
圈数取整数，故绳子至少需绕 \(\boldsymbol{15}\) 圈。
\\
\\
\\
像我这样不懂什么欧拉公式的同学，可以像下面这样做\\
取绳上一微元段，对应圆心角 $d\theta$，两端张力分别为 $T$ 与 $T+dT$。
微元受到圆柱法向支持力 $dN$ 与静摩擦力 $df=\mu dN$。
切向（沿切线方向）平衡：
\[
(T+dT)\cos\frac{d\theta}{2}
= T\cos\frac{d\theta}{2} + \mu dN
\]
法向（沿半径方向）平衡：
\[
dN
= (T+dT)\sin\frac{d\theta}{2} + T\sin\frac{d\theta}{2}
\]
小角度近似：
$\cos\frac{d\theta}{2}\approx1,\ \sin\frac{d\theta}{2}\approx\frac{d\theta}{2}$，
略去高阶小量 $dT\cdot d\theta$，化简得：
\[
dT = \mu T d\theta
\]
设轻端张力 $T_1=mg$，重端张力 $T_2=m_0g$，
总包角为 $\theta$，分离变量积分：
\[
\int_{T_1}^{T_2}\frac{dT}{T}
= \mu\int_0^\theta d\theta
\]
\[
\ln\frac{T_2}{T_1} = \mu\theta
\]
代入 $T_2=m_0g,\ T_1=mg$，约去 $g$：
\[
\ln\frac{m_0}{m} = \mu\theta
\]
设圈数为 $n$，则 $\theta=2\pi n$，故
\[
n = \frac{\ln(m_0/m)}{2\pi\mu}
\]
代入数据 $m_0=1000\ \text{kg},\ m=10\ \text{kg},\ \mu=0.05$：
\[
n=\frac{\ln 100}{2\pi\times 0.05}\approx 14.66
\]
圈数取整数，至少需绕 $\boldsymbol{15}$ 圈。
    }
\end{enumerate}

\section{三、动量}
\begin{enumerate}[label=\arabic*.,leftmargin=*]
    \item 冲量是力对时间的积累。怎么理解这句话？\\
    \textcolor{blue}{就是在时间上的积累，推箱子不动但动量不变是因为动量定理看的是合外力，此时箱子所受合外力为零}

    \item 阅读3.1，合力的冲量等于各个分力冲量的矢量和。那么合力的做功是否等于分力做功的和呢？试图证明或是证伪。\\
    \textcolor{blue}{由于合力为各力的矢量和：
\[
\boldsymbol{F}_{\text{合}}=\sum_{i=1}^{n}\boldsymbol{F}_{i}
\]
故合力做功
\[
W=\int \boldsymbol{F}_{\text{合}}\cdot d\boldsymbol{s}
=\int \sum_{i=1}^{n}\boldsymbol{F}_{i}\cdot d\boldsymbol{s}
=\sum_{i=1}^{n}\int \boldsymbol{F}_{i}\cdot d\boldsymbol{s}
\]
（$\boldsymbol{F}$ 和 $d\boldsymbol{s}$ 均为矢量，中间为点乘）证毕。}

    \item 阅读3.2.2“在应用动量守恒定律时应注意以下几点：”，
    
    （3）在质点斜碰或者球非对心碰撞（就是初速度方向与两球心连线方向不共线，如例3-5）的时候会有所运用；
    
    \colorbox{yellow}{括号里的描述似乎是错误的，“初速度方向与两球心连线方向不共线”}

    （4）在不涉及转动的时候，由于牛顿定律有时候会比较麻烦，所以可以用动量约束和能量约束列方程组，从而舍弃牛顿定律的复杂运算。除非题目让你算加速度。

    \item 阅读例3-3.你发现同时跳和依次跳的末速度是不同的。为什么？如果我有100个人又应该得到什么结果呢？\\
    \colorbox{yellow}{"一起跳"应更正为"依次跳"}

    \textcolor{blue}{本质上是因为跳下时的速度是相对速度，分次跳下相比于同时跳下，可以感觉到人每次跳下的对地速度是在增大的，故车的末动量更大}

    \item 例3-4，本题在先修考试中有涉及。\\
    \textcolor{blue}{ft=mv}

    \item 阅读3.5，观察式3-25，对该式两端对时间微分和积分，得到两个新式子。再结合式3-26与例3-9，希望对你有所启发。

    \colorbox{yellow}{式3-25两端对时间微分得到质心加速度与质心外力的关系；积分得到什么？}

    \textcolor{blue}{积分会得到一维的质心坐标的式子}

    \textcolor{blue}{质心式子的推广，关键是看里面那个物理量是相对于谁而言的，有时可以得出“相对质心的质心运动速度”，即为0}
    \item （不会考，“邪修”方法，可以不看）阅读例3-8，搜索“巴普斯定理”。
    
    \textcolor{blue}{\textbf{巴普斯定理：} 平面均质物体垂直平面运动扫过的体积 = 物体面积 × 质心运动路程
\\例题：半径$R$的均质半圆盘质心位置
\[
\frac{4}{3}\pi R^3 = \frac{1}{2}\pi R^2 \cdot 2\pi x
\implies x = \frac{4R}{3\pi}
\]
\textbf{均质细铁丝半圆环质心（定义式计算）}
设半圆环半径$R$，总质量$m$，线密度$\lambda = \dfrac{m}{\pi R}$，极角$\theta\in[0,\pi]$，坐标$x=R\cos\theta,y=R\sin\theta$，质量元$dm=\lambda R d\theta$。
质心定义：
\[
x_c = \frac{1}{m}\int x dm, \quad y_c = \frac{1}{m}\int y dm
\]
$x$方向（对称抵消）：
\[
x_c = \frac{1}{m}\int_0^\pi R\cos\theta·\lambda R d\theta = \frac{\lambda R^2}{m}\int_0^\pi \cos\theta d\theta = 0
\]
$y$方向：
\begin{align*}
y_c &= \frac{1}{m}\int_0^\pi R\sin\theta·\lambda R d\theta = \frac{\lambda R^2}{m}\int_0^\pi \sin\theta d\theta \\
&= \frac{\lambda R^2}{m}·2 = \frac{2·\frac{m}{\pi R}·R^2}{m} = \frac{2R}{\pi}
\end{align*}
结论：质心在对称轴上，距圆心$\boldsymbol{\dfrac{2R}{\pi}}$。
    }

    \item 习题中需要注意的地方：3-7,3-8（巩固对于例3-3的记忆，并且了解瞬间发射和均匀喷出的区别）；3-9（注意二人质量相等的结果）；
    3-12（思考当我的秤的读数为某一个值的时候，我应该在哪里截断才能保证我的称出的重量是准确的？）；3-13（难题，考验计算）\\
    \textcolor{blue}{
    3-7瞬间喷气：只受一次冲量，过程近似动量守恒\\
    均匀喷气：要用神秘变质量物体运动方程（密歇尔斯基）\\
    3-9
两个质量分别为 $m$ 和 $m_0$ ($m_0 > m$) 的人分别挂在定滑轮两侧的绳子上。设绳子不可伸长，滑轮及绳子的质量、滑轮轴承处的摩擦均可忽略不计，滑轮两边绳子的张力大小相同。若开始时，两人离滑轮的距离都是 $h$。试证明：当质量较轻的人在 $t$ 时刻爬到滑轮处时，质量较重的人与滑轮的距离为
\[
\frac{m_0 - m}{m_0}\left(h + \frac{1}{2}gt^2\right)
\]
注意二者的加速度方向不一致（列式时可以假定a加速度方向，算出正负代表方向），所以消元掉T时要注意符号。\\
后续积分求解
\begin{align*}
\text{对轻人 (质量 } m\text{): } & F_T - mg = ma_1 = m\frac{d^2x_1}{dt^2} \\
\text{对重人 (质量 } m_0\text{): } & F_T - m_0g = m_0a_2 = m_0\frac{d^2x_2}{dt^2}
\end{align*}
\textbf{2. 加速度关系}
两式相减消去 $F_T$，得到加速度关系：
\[
(m_0 - m)g = m\frac{d^2x_1}{dt^2} - m_0\frac{d^2x_2}{dt^2}
\]
\textbf{3. 积分求解速度}
对时间 $t$ 进行第一次积分（初始速度均为 0）：
\[
\int_0^t (m_0 - m)g \, dt = m\int_0^{v_1} dv_1 - m_0\int_0^{v_2} dv_2
\]
\[
(m_0 - m)gt = mv_{1} - m_0v_{2}
\]
即：
\[
(m_0 - m)gt = m\frac{dx_1}{dt} - m_0\frac{dx_2}{dt}
\]
\textbf{4. 积分求解位移}
再次对时间 $t$ 积分（初始位移均为 $h$）：
\[
\int_0^t (m_0 - m)gt \, dt = m\int_{-h}^{x_1} dx_1 - m_0\int_{-h}^{x_2} dx_2
\]
注：轻人向上运动，位移变化为 $x_1 - (-h)$；重人位移变化为 $x_2 - (-h)$。
\[
\frac{1}{2}(m_0 - m)gt^2 = m(x_1 + h) - m_0(x_2 + h)
\]
\textbf{5. 代入边界条件求距离}
题目条件：轻人到达滑轮处，即 $x_1 = 0$。
代入上式：
\[
\frac{1}{2}(m_0 - m)gt^2 = m(0 + h) - m_0(x_2 + h)
\]
整理求解重人到滑轮的距离 $l = |x_2|$（因 $x_2$ 为负，代表在滑轮下方）：
\[
m_0(x_2 + h) = mh - \frac{1}{2}(m_0 - m)gt^2
\]
\[
x_2 = \frac{mh}{m_0} + \frac{m - m_0}{2m_0}gt^2 - h
\]
\[
x_2 = - \frac{m_0 - m}{m_0}h - \frac{m_0 - m}{2m_0}gt^2
\]
因此，重人与滑轮的距离 $l$ 为：
\[
l = \left| x_2 \right| = \frac{m_0 - m}{m_0} \left( h + \frac{1}{2}gt^2 \right)
\]
\textbf{证毕。}
\\3-12流体问题可以不考虑到底有多少还留在空中，所以就是直接按最简单的想法列式就好
\\3-13半球也会运动哦,用半球参考系\\
\textbf{1. 水平方向动量守恒}
系统水平方向不受外力，动量守恒。设半球对地速度为$V$（向右为正），滑块相对半球的速度为$v'$，则滑块对地的水平速度为$v'\cos\theta - V$。
\[
MV = m(v'\cos\theta - V)
\]
整理得：
\[
(M + m)V = mv'\cos\theta \tag{1}
\]
\textbf{2. 系统机械能守恒}
半球光滑、水平面光滑，系统机械能守恒。滑块下落高度$h = R(1-\cos\theta)$，重力势能转化为滑块与半球的动能：
\[
mgR(1-\cos\theta) = \frac{1}{2}m\left[(v'\cos\theta - V)^2 + (v'\sin\theta)^2\right] + \frac{1}{2}MV^2
\]
展开滑块速度平方项：
\[
(v'\cos\theta - V)^2 + (v'\sin\theta)^2 = v'^2 - 2v'V\cos\theta + V^2
\]
代入得：
\[
mgR(1-\cos\theta) = \frac{1}{2}m(v'^2 - 2v'V\cos\theta + V^2) + \frac{1}{2}MV^2
\]
\textbf{3. 联立求解$v'$}
将式(1) $V = \dfrac{mv'\cos\theta}{M+m}$ 代入上式，化简后解得：
\[
v' = \left[ \frac{2gR(1-\cos\theta)(m+M)}{M + m\sin^2\theta} \right]^{\frac{1}{2}}
\]
\textbf{4. 脱离条件（法向加速度）}
滑块脱离半球瞬间，支持力$N=0$，法向合力提供向心力（以半球为参考系，相对加速度）：
\[
g\cos\theta = \frac{v'^2}{R}
\]
\textbf{5. 代入$v'$求解质量比}
将$v'^2 = \dfrac{2gR(1-\cos\theta)(M+m)}{M + m\sin^2\theta}$ 代入上式，约去$gR$：
\[
\cos\theta = \frac{2(1-\cos\theta)(M+m)}{M + m\sin^2\theta}
\]
利用$\sin^2\theta = 1-\cos^2\theta$，整理得：
\[
\frac{M}{m} = \frac{2(1-\cos\theta) - (1-\cos^2\theta)\cos\theta}{\cos\theta - 2(1-\cos\theta)}
\]
代入$\cos\theta = 0.7$：
\begin{align*}
\frac{M}{m} &= \frac{2(1-0.7) - (1-0.7^2)\times0.7}{0.7 - 2(1-0.7)} \\
&= \frac{2\times0.3 - (1-0.49)\times0.7}{0.7 - 2\times0.3} \\
&= \frac{0.6 - 0.51\times0.7}{0.7 - 0.6} \\
&= \frac{0.6 - 0.357}{0.1} = \frac{0.243}{0.1} = 2.43
\end{align*}
    }
\end{enumerate}

\section{四、能量}
\begin{enumerate}[label=\arabic*.,leftmargin=*]
    \item 阅读4.2。思考，为什么摩擦力不是保守力？如果我有一个力，其表达式可以写成
    \[
    \boldsymbol{F}=kx\boldsymbol{i}
    \]
    那么这个力是保守力吗？进一步的，如果这个力是
    \[
    \boldsymbol{F}=f(x)\boldsymbol{i}+g(y)\boldsymbol{j}+w(z)\boldsymbol{k}
    \]
    情况又如何呢？另外，尝试证明弹簧弹力是保守力。\\
    \textcolor{blue}{
    做功只与物体起点和终点位置有关，而与路径无关的力是保守力；\\    
    摩擦力的做功与路径有关；\\
    如果力的表达式可以写成某个函数的梯度（即存在势能函数），则该力是保守力；\\
    弹簧弹力满足$F=-kx$，势能函数为$U=\frac{1}{2}kx^2$，满足保守力定义。\\
    \textbf{2. 力是梯度 $\implies$ 存在势能函数}
若力 $\vec{F}(\vec{r})$ 可表示为某个标量函数 $U(\vec{r})$ 的梯度：
\begin{align*}
\vec{F} = -\nabla U
\end{align*}
则力对任意路径 $L$（从 $\vec{r}_A$ 到 $\vec{r}_B$）的做功为：
\begin{align*}
W_{AB} &= \int_{\vec{r}_A}^{\vec{r}_B} \vec{F} \cdot d\vec{r} = \int_{\vec{r}_A}^{\vec{r}_B} (-\nabla U) \cdot d\vec{r}
\end{align*}
由多元函数全微分公式 $dU = \nabla U \cdot d\vec{r}$，代入得：
\begin{align*}
W_{AB} = -\int_{\vec{r}_A}^{\vec{r}_B} dU = U(\vec{r}_A) - U(\vec{r}_B)
\end{align*}
做功仅由初末位置的 $U$ 值决定，与路径无关，因此 $\vec{F}$ 是保守力，$U(\vec{r})$ 即为对应的势能函数。
    }

    \item 返回去阅读例4-4，求解力的平衡位置。这是否能引发你的思考？（势能？）
\textcolor{blue}{\\平衡位置势能为极值\\
\textbf{三种平衡状态的对应}
\begin{align*}
& \text{1. 稳定平衡：} & U(\vec{r}) \text{ 取极小值} \\
& \text{2. 不稳定平衡：} & U(\vec{r}) \text{ 取极大值} \\
& \text{3. 随遇平衡：} & U(\vec{r}) \text{ 取常值（无极值）}
\end{align*}}

    \item 保守力、势能、场知一推二。可以自己试图证明。势，potential，代表了保守力做功的潜力。一般机械能就像水一样，重力势能是一个水箱，弹性势能是一个水箱。平时这两个水箱里的水都是藏起来而不示人的，而动能就像是接出来在水杯里看得见的水。各非保守力消耗的机械能就好像水撒在了地上。
\textcolor{blue}{\\对}
    \item 柯尼希定理的证明。自己证明一遍，数学上保持严谨。
\textcolor{blue}{\\也许就是相对运动公式}
    \item 阅读4.3.3，“对于多个物体组成的系统而言，外力和非保守内力做的功之和等于系统机械能的增量”。谈谈你的理解。
\textcolor{blue}{\\外力如推力，非保守内力如爆炸，他们做功都会机械能}
    \item 阅读例4-8，尝试用牛顿定律解决。感受下动量与能量的优越性。
\textcolor{blue}{\\啊对}
    \item 习题：试图证明，对于光滑与地面不固结的轨道，有一个小球以某一初速度滑上轨道，且不会脱离。当小球处于最高点时，小球与轨道共速，即使轨道形状千奇百怪。并求出小球的最高高度。
\textcolor{blue}{\\不共速意味着还能上升，即不是最高点，最高高度用动量加能量}
    \item 习题中需要注意的地方：4-14（寻找积分方法是比较重要的。当然，可以参考邪修方法巴普斯定理）；4-17（困难）；4-19（困难。《难题集萃》研究生物理模拟试题（一）第一题）。
\textcolor{blue}{\\4-14直接积分就好，设面密度
\\4-17和4-19暂时没空看}
\end{enumerate}

\section{五、角动量}
\begin{enumerate}[label=\arabic*.,leftmargin=*]
    \item 阅读例5-2.了解到开普勒第二定律的实质就是角动量守恒定律。也就是说，对于某一天体的有心运动来说，角动量是守恒的。即，只要受力方向与参考点到该点的方向相同，则角动量守恒。包括力为0和距离为0的情况。
    \textcolor{blue}{\\椭圆，面积，调和}\\
    \colorbox{yellow}{"距离"改成“力臂”是否更恰当}
    \textcolor{blue}{\\力为零，如匀速直线运动，对任意固定参考点，角动量不变：
\begin{align*}
\vec{L} &= \vec{r} \times m\vec{v} = \text{常矢量}
\end{align*}}
    \item 阅读5.2，注意到式5-10.随后思考本式子与平行轴定理的联系。
\textcolor{blue}{\\兴许是有的}
    \item （夹带私货）《题金山寺》真乃好诗也！
\textcolor{blue}{\\遥望四边云接水，碧峰千点数鸿轻。
\\轻鸿数点千峰碧，水接云边四望遥。
\\P107}
    \item 可以阅读一下5.3.3和5.3.4，还是很有意思的。
\textcolor{blue}{\\对称性}
    \item 习题需要注意的地方：5-6（注意某些特殊情况）；5-8（用数学解释，也可以抽象地用对称性解释）；5-9（角动量守恒吗？关于哪一点呢？）。
\textcolor{blue}{\\5-6最开始的时候，列力矩的式子（F和G）
\\后面的题目以后再来看}
\end{enumerate}

\section{六、刚体力学}
\begin{enumerate}[label=\arabic*.,leftmargin=*]
    \item 回答章前问题。
\textcolor{blue}{\\减小转动惯量}
    \item 阅读6.2.1，如果合外力力矩方向和轴的方向不相同会发生什么？如果刚体并不是定轴转动而是定点转动又会发生什么？
\textcolor{blue}{\\不沿轴的分量试图让轴倾斜，但是定轴运动，所以产生约束力矩，依然绕原轴转动
\\定点转动：\\
- 力矩与角动量同向：角动量大小变化，转轴方向不变
\\- 力矩与角动量垂直：角动量方向旋转，产生进动
\\- 力矩任意方向：角动量大小与方向同时变化，出现进动与章动(物体旋转发生进动时，转轴R在绕L转动的过程中发生轻微抖动（铅垂面内上下摆动）的现象。)}
    \item 阅读表6-1.全部进行证明。另外，可以试图证明均匀圆盘对某一半径的转动惯量是多少。
\textcolor{blue}{\\平行轴}
    \item 阅读6.2.3，与动量守恒定律进行对比研究。同样，转动动能也和平动动能有着数学上的相似性，是物理独特的美感。

    \item 阅读例6-5，尝试自己证明一遍.
\textcolor{blue}{\\关键是取烟囱截面分析，考虑惯性力，求导分析力矩最大的地方}
    \item 习题中需要注意的地方：习题6-6（打击中心的概念）；习题6-9（你就说吊着是不是相对于绳匀速吧）；习题6-12（得出具体数值解。现在我有一个质量均匀分布的圆盘也是这样转，求转动惯量和摩擦力矩。为什么会出现这样的答案呢？）
\textcolor{blue}{\\6-9好邪门的想法，老实做的话是利用加速度相等
\\6-12好题}
\end{enumerate}

\section{七、流体力学}
\begin{enumerate}[label=\arabic*.,leftmargin=*]
    \item 阅读7.2，想一想，“单位时间内流入任意闭合曲面的流体的体积一定与流出该闭合曲面的流体的体积相等”。这是一个重要的约束。想一想，基尔霍夫电流定律是不是也差不多？如果你有相关基础，是否可以发现流体连续性原理、基尔霍夫电流定律、电场高斯定律之间共有的数学关系？
\textcolor{blue}{\\无源场，流入 = 流出，
闭合面上的总通量 = 0}
    \item 请证明伯努利方程。（先修考试第一道大题）
\textcolor{blue}{\\面积功加体积功，功能关系/动能定理，连续性方程->可解}
    \item 请解释虹吸效应。既然虹吸效应可以把水柱提高一定的高度，为什么还会有耗电的抽水机存在呢？
\textcolor{blue}{\\虹吸效应是由于大气压强和重力作用形成的，水柱高度受限于大气压强，最多上升10m，要求出口比入口低；抽水机可以提供更大的压力差，克服更高的水柱高度限制，且可以精准控制。}
\end{enumerate}

\end{document}
